% Standalone section source for the parent article. Uses standard amsmath/amsthm.
\section{Coincident reflected Stieltjes products close in ordinary zeta jets}
\label{ct:section}

Let
\[
 K(x)=\pi\cot(\pi x),\qquad
 \zeta(1+u,x)-\frac1u
   =\sum_{m\ge0}\frac{(-1)^m}{m!}\gamma_m(x)u^m.
\]
In particular $\gamma_0(x)=-\psi(x)$. All finite parts in this section use
\emph{both} unit endpoint coordinates $x$ and $1-x$. Equivalently,
\[
 \operatorname{FP}\int_0^1 f(x)\,dx
 =\operatorname{CT}_{\epsilon\downarrow0}
       \int_\epsilon^{1-\epsilon} f(x)\,dx,
\]
where the constant coefficient is taken after expanding every endpoint
power and logarithm. Independent lower and upper cutoffs give the same
answer because the two local contributions are added. This convention
is invariant under $x\mapsto1-x$, and it assigns zero to
$\operatorname{FP}\int_0^1x^{-1}\log^j x\,dx$.

The incoming triple report proves finite closure for one Stieltjes factor
\cite{RepoTriple} and any number of reflected polygamma factors whose seams are pairwise
separated. The result below allows every reflected factor to collide with
the Stieltjes seam. It is an exact closed subclass of the coincident
product problem. It does not evaluate the unreflected cube
$\operatorname{FP}\int\psi^3$, whose remaining diagonal Tornheim third
jet is a different coordinate. The Fourier transform and reflection
inputs are classical; the contribution here is the full coordinate
completion and its all-order coincident reflected consequences.

\subsection{A classical meromorphic transform and its local completion}

\begin{lemma}[Cotangent--Hurwitz transform]
\label{ct:transform}
For $r\ge0$, initially when $\Re s$ is sufficiently negative and then by
meromorphic continuation,
\begin{equation}\label{ct:J}
 J_r(s):=\operatorname{FP}\int_0^1\zeta(s,x)K^{(r)}(x)\,dx
 =(s)_r\,\pi\cot\frac{\pi(s+r)}2\,\zeta(s+r).
\end{equation}
At a spectral resonance the right side denotes the meromorphic
continuation, not the coordinate finite part of a specialized
coefficient.
\end{lemma}
\begin{proof}
The periodic principal value of $K$ has nonzero Fourier coefficients
$-\pi i\operatorname{sgn}(n)$ and zero mean. This follows, for example,
by differentiating the Abel Fourier series for
$\log(2\sin\pi x)$ in distributions. The Hurwitz Fourier expansion \cite[Eq.~25.11.9]{DLMFHurwitz} gives,
for $\Re s<0$,
\[
 J_0(s)=2\pi\Gamma(1-s)(2\pi)^{s-1}
             \cos(\pi s/2)\zeta(1-s).
\]
The sum is absolutely convergent because its summands have magnitude
$O(n^{\Re s-1})$. The Riemann functional equation \cite{DLMFZetaReflection} turns the displayed
expression into $\pi\cot(\pi s/2)\zeta(s)$, with removable singularities
interpreted by continuation.

For $\Re s<-r-1$, the singular term $x^{-s}$ and sufficiently many
of its derivatives vanish at zero. The smooth part
$\zeta(s,1+x)$ has the same periodic endpoint jets as $\zeta(s,x)$ at
one. Consequently distributional integration by parts against
$K^{(r)}$ agrees with the stated endpoint finite part and gives
\[
 J_r(s)=(-1)^r\operatorname{FP}\int_0^1
              \partial_x^r\zeta(s,x)K(x)\,dx
       =(s)_rJ_0(s+r).
\]
There are no unrecorded endpoint constants in this initial region:
the smooth jets match and the remaining power vanishes. Finite endpoint
Taylor subtraction supplies meromorphic continuation in $s$, proving
\eqref{ct:J} everywhere as a meromorphic identity.
\end{proof}

Set $N=p+r$ and define $P_j(u)=(1+u)_j$, so $P_0=1$.
For $p,r\ge0$ put
\begin{align}
 \mathcal R_{p,r}(u)
 =(-1)^p\bigg[&P_N(u)\,
      \pi\cot\frac{\pi(N+1+u)}2\,\zeta(N+1+u)\notag\\
 &-\mathbf1_{N\text{ odd}}
      \frac{2N!\zeta(N+1)}{p!}\,\frac{P_p(u)}u\bigg].
 \label{ct:R}
\end{align}
For $N=0$ the first line is the removable expression
$-\pi\tan(\pi u/2)\zeta(1+u)$.

\begin{theorem}[All-index coincident reflected closure]
\label{ct:closure}
The function $\mathcal R_{p,r}$ is holomorphic near $u=0$, and for every
$m,p,r\ge0$,
\begin{equation}\label{ct:moment}
 \boxed{\displaystyle
 \operatorname{FP}\int_0^1\gamma_m^{(p)}(x)K^{(r)}(x)\,dx
       =(-1)^m m!\,[u^m]\mathcal R_{p,r}(u).}
\end{equation}
Thus all these moments have a finite reduction to ordinary Riemann zeta
jets at $N+1$ and generalized harmonic polynomials. In the exceptional
case $N=0$, ordinary Stieltjes constants occur instead.
\end{theorem}
\begin{proof}
Write $g(u,x)=\zeta(1+u,x)-1/u$. The exact local splitting is
\[
 \partial_x^pg(u,x)
   =(-1)^pP_p(u)x^{-p-1-u}+h_p(u,x),\qquad
 h_p(u,x)=\partial_x^pg(u,1+x).
\]
The function $h_p$ is holomorphic jointly near $(u,x)=(0,0)$.
At the upper endpoint the entire factor $\partial_x^pg(u,1-y)$ is
holomorphic jointly near $(0,0)$ in $(u,y)$. The cotangent expansion is
\[
 K^{(r)}(x)=(-1)^rr!x^{-r-1}
 -2\sum_{j\ge1}\zeta(2j)
        \frac{(2j-1)!}{(2j-1-r)!}x^{2j-1-r},
\]
where a term is zero when $2j-1<r$.
In the product of the varying singular power with this expansion,
resonance with $x^{-1-u}$ occurs exactly when
$2j=N+1$, hence exactly when $N$ is odd. Its complete coefficient is
\[
 A_{p,r}(u)
 =-2(-1)^p\frac{N!}{p!}\zeta(N+1)P_p(u).
\]
The meromorphic integral of this local monomial on $(0,b)$ is
$-A_{p,r}(u)b^{-u}/u$. Its coefficientwise coordinate finite-part
generator is $-A_{p,r}(u)(b^{-u}-1)/u$. Their difference is
$-A_{p,r}(u)/u$; therefore the required completion is the meromorphic
integral plus $A_{p,r}(u)/u$.

All other varying-power monomials have denominators nonzero at zero.
Subtract finitely many local Taylor terms at both endpoints so that all
remainders and any fixed finite number of parameter derivatives are
bounded by $Cx^{-\eta}(1+|\log x|^M)$, respectively
$Cy^{-\eta}(1+|\log y|^M)$, with $\eta<1$. These integrals are normally
convergent. The fixed integer endpoint powers from the smooth factors
have analytic coefficients and their coordinate finite parts may be
extracted coefficientwise. This proves holomorphy and identifies all
Taylor coefficients.

The meromorphic integral itself is, by \eqref{ct:J},
\[
 (-1)^pP_p(u)J_r(1+p+u)
 =(-1)^pP_N(u)\pi\cot\frac{\pi(N+1+u)}2\zeta(N+1+u).
\]
The subtraction $-K^{(r)}/u$ when $p=0$ contributes nothing because
$K^{(r)}$ has coordinate finite-part mean zero. Adding the complete
local correction just obtained gives exactly \eqref{ct:R}, and
coefficient extraction in $g$ proves \eqref{ct:moment}.
\end{proof}

Removing only the pole residue would replace $P_p(u)$ by $p!$ and
lose finite terms. In particular, for odd $N$ it would incorrectly
omit the harmonic correction $H_p$ in the next formula.

\subsection{Polygamma specializations and exact finite endpoint anomalies}

Let $H_0=0$. At Stieltjes index zero, \eqref{ct:moment} gives
\begin{equation}\label{ct:zero-index}
 \operatorname{FP}\int_0^1\psi^{(p)}(x)K^{(r)}(x)\,dx
 =\begin{cases}
 \pi^2/2,&N=0,\\
 0,&N>0\text{ even},\\
 2(-1)^{p+1}N!\bigl[\zeta'(N+1)
                +(H_N-H_p)\zeta(N+1)\bigr],&N\text{ odd}.
 \end{cases}
\end{equation}
The first examples are
\begin{align*}
 \operatorname{FP}\int_0^1\psi(x)K'(x)\,dx
   &=-2\zeta'(2)-2\zeta(2),\\
 \operatorname{FP}\int_0^1\psi'(x)K(x)\,dx
   &=2\zeta'(2),\\
 \operatorname{FP}\int_0^1\psi(x)K'''(x)\,dx
   &=-12\zeta'(4)-22\zeta(4),\\
 \operatorname{FP}\int_0^1\psi'(x)K''(x)\,dx
   &=12\zeta'(4)+10\zeta(4),\\
 \operatorname{FP}\int_0^1\psi''(x)K'(x)\,dx
   &=-12\zeta'(4)-4\zeta(4),\\
 \operatorname{FP}\int_0^1\psi'''(x)K(x)\,dx
   &=12\zeta'(4).
\end{align*}
These are consistent with finite-part integration by parts, including
its boundary constants. For example,
\[
 \operatorname{FP}\int_0^1(\psi K)'\,dx=-2\zeta(2).
\]
Indeed the constant terms of $\psi(x)K(x)$ at zero and one are
$3\zeta(2)$ and $\zeta(2)$, respectively. Thus setting the integral
of this ordinary derivative to zero would be incorrect in the present
coordinate convention.

At $p=r=0$, every Stieltjes index has the particularly short formula
\begin{align}
 \operatorname{FP}\int_0^1\gamma_m(x)K(x)\,dx
 ={}&4m!\sum_{k=1}^{\lfloor(m+1)/2\rfloor}
   \frac{(1-2^{-2k})\zeta(2k)\gamma_{m-2k+1}}
                       {(m-2k+1)!}\notag\\
 &-\mathbf1_{2\mid m}\,4m!(1-2^{-m-2})\zeta(m+2).
 \label{ct:all-m-base}
\end{align}
For $m=0,1,2$ this gives
\[
 -\frac{\pi^2}{2},\qquad
 \frac{\pi^2\gamma}{2},\qquad
 \pi^2\gamma_1-\frac{\pi^4}{12}.
\]
The base constant can also be recovered formally from the finite
constant of the separated Stieltjes--cotangent transform in the incoming
triple report. The local proof above establishes that this formal limit
has the stated same-point coordinate normalization, and supplies all
argument-derivative resonances.

For $N>0$ even and $m=1$ there is the simple additional family
\[
 \operatorname{FP}\int_0^1\gamma_1^{(p)}(x)K^{(r)}(x)\,dx
   =\frac{(-1)^pN!\pi^2}{2}\zeta(N+1).
\]
All higher coefficients in \eqref{ct:R} are finite because
$P_j(u)=j!\prod_{a=1}^j(1+u/a)$ is a polynomial.

\subsection{Arbitrarily many coincident reflected factors}

Define derivative polynomials
\[
 D_0(X)=X,\qquad D_{r+1}(X)=-(X^2+\pi^2)D_r'(X).
\]
Then $K^{(r)}(x)=D_r(K(x))$, and $D_r$ has degree $r+1$ and leading
coefficient $(-1)^rr!$. Consequently every polynomial $P(X)$ has a unique
finite decomposition
\begin{equation}\label{ct:poly-decomp}
 P(X)=C+\sum_{r=0}^{\deg P-1}c_rD_r(X).
\end{equation}
The constant is explicitly
\[
 C=\frac{P(i\pi)+P(-i\pi)}2,
\]
since $D_r(\pm i\pi)=0$ for $r\ge1$ and $D_0(\pm i\pi)=\pm i\pi$.
The remaining coefficients follow by descending elimination of the
leading powers, an exact terminating procedure.

\begin{corollary}[Finite coincident reflected algebra]
\label{ct:algebra}
For arbitrary $m,p\ge0$ and any polynomial $P$, the finite part
\[
 \operatorname{FP}\int_0^1\gamma_m^{(p)}(x)P(K(x))\,dx
 =\sum_r c_r(-1)^m m![u^m]\mathcal R_{p,r}(u)
\]
has a finite evaluation in ordinary zeta jets and Stieltjes constants.
The same holds after replacing
$P(K)$ by any finite product of arbitrary derivatives $K^{(r_j)}$.
Equivalently, it holds for any finite product of the reflected factors
\[
 (-1)^{r_j}\psi^{(r_j)}(1-x)-\psi^{(r_j)}(x).
\]
\end{corollary}
\begin{proof}
The reflection formula \cite[Eq.~5.15.6]{DLMFPolyGamma} identifies each balanced polygamma expression
with $K^{(r_j)}$. The derivative-polynomial recurrence converts their
product to a polynomial in $K$. Apply \eqref{ct:poly-decomp} and
Theorem~\ref{ct:closure}. The constant term contributes zero because
$\operatorname{FP}\int_0^1\gamma_m^{(p)}(x)\,dx=0$. For $p=0$ this
follows by integrating the shifted Hurwitz family and completing its
single resonance; for $p\ge1$ it follows directly by taking endpoint
constant terms of $\gamma_m^{(p-1)}$. The smooth endpoint jets coincide
and the singular derivative of $x^{-1}\log^m x$ has no constant term.
\end{proof}

For powers alone, write
$K^q=C_q+\sum_rc_{q,r}K^{(r)}$. A useful recurrence is
\[
 C_0=1,\quad C_1=0,\quad c_{1,0}=1,\qquad
 C_q=-\pi^2 C_{q-2},\qquad
 c_{q,r}=-\frac{c_{q-1,r-1}}{q-1}-\pi^2c_{q-2,r},
\]
with absent coefficients zero. It follows from
$(K^{q-1})'=-(q-1)(K^q+\pi^2K^{q-2})$.
In particular,
\begin{align*}
 K^2&=-K'-\pi^2,\\
 K^3&=\tfrac12K''-\pi^2K,\\
 K^4&=-\tfrac16K'''+\tfrac43\pi^2K'+\pi^4,\\
 K^5&=\tfrac1{24}K^{(4)}-\tfrac56\pi^2K''+\pi^4K.
\end{align*}
Since the coordinate finite-part mean of every $K^{(r)}$ is zero,
\begin{equation}\label{ct:pure-power}
 \operatorname{FP}\int_0^1K(x)^q\,dx
       =\pi^q\cos(q\pi/2).
\end{equation}
The resulting mixed digamma examples include
\begin{align}
 \operatorname{FP}\int_0^1\psi(x)K(x)^2\,dx
     &=2\zeta'(2)+2\zeta(2),\label{ct:cubic-reflected}\\
 \operatorname{FP}\int_0^1\psi(x)K(x)^4\,dx
     &=2\zeta'(4)-\frac{8\pi^2}{3}\zeta'(2)
                       -\frac{109}{3}\zeta(4).\label{ct:quintic-reflected}
\end{align}
There is an all-degree elementary family:
\begin{equation}\label{ct:odd-powers}
 \boxed{\displaystyle
 \operatorname{FP}\int_0^1\psi(x)K(x)^{2h+1}\,dx
       =\frac{(-1)^h\pi^{2h+2}}2\qquad(h\ge0).}
\end{equation}
Besides the polynomial proof, reflection gives a short independent
argument: averaging the integral with its reflected copy yields
$-\tfrac12\operatorname{FP}\int K^{2h+2}$, which is
\eqref{ct:odd-powers} by \eqref{ct:pure-power}.

There is also an exactly reducible family with two unreflected factors:
\begin{equation}\label{ct:two-unreflected}
 \operatorname{FP}\int_0^1\psi(x)^2K(x)^{2h+1}\,dx
   =-\operatorname{FP}\int_0^1\psi(x)K(x)^{2h+2}\,dx.
\end{equation}
To prove it, reflect the left side and use
$\psi(1-x)=\psi(x)+K(x)$. The reflected value is its negative minus
twice the right-hand moment before its displayed minus sign, while
the remaining pure odd power has zero finite-part mean. Thus every
moment in \eqref{ct:two-unreflected} also has a finite ordinary-zeta-jet
reduction by Corollary~\ref{ct:algebra}. At $h=0$ it says
\[
 \operatorname{FP}\int_0^1\psi(x)^2K(x)\,dx
       =-2\zeta'(2)-2\zeta(2).
\]

A direct cubic consequence, retaining the distinction from the open
unreflected cube, is
\begin{equation}\label{ct:cubic-difference}
 \operatorname{FP}\int_0^1
 \bigl[\psi(x)^3-\psi(x)^2\psi(1-x)\bigr]dx
      =2\zeta'(2)+2\zeta(2).
\end{equation}
Indeed expanding $\psi K^2$ gives
$\psi\psi(1-x)^2-2\psi^2\psi(1-x)+\psi^3$, and reflection identifies
the finite parts of the first two mixed cubic monomials. This determines
a difference of cubic moments and leaves their common unreflected
Tornheim coordinate free.

\subsection{Log-Gamma primitives and a harmonic-number family}

The same completed transform supplies the integrated Gamma level.
At $s=0$, the only resonance for $\zeta(s,x)K(x)$ is $x^{-1-s}$.
Thus
\[
 \pi\cot(\pi s/2)\zeta(s)+\frac1s
\]
is the generating germ for the coordinate finite parts of the spectral
jets $\zeta^{[j]}(0,x)K(x)$. Its linear coefficient is
$\zeta''(0)+\zeta(2)/2$, because
$\zeta(0)=-1/2$ and
$\pi\cot(\pi s/2)=2/s-\zeta(2)s+O(s^3)$.
Lerch's identity \cite[Eq.~25.11.18]{DLMFHurwitz} $\zeta'(0,x)=\log\Gamma(x)-\tfrac12\log(2\pi)$
and the zero mean of $K$ give
\begin{equation}\label{ct:loggamma-K}
 \operatorname{FP}\int_0^1\log\Gamma(x)K(x)\,dx
       =\zeta''(0)+\frac12\zeta(2).
\end{equation}
The constant $1/s$ is independent of the parameter, apart from its
pole, so this specialization has no omitted finite derivative of a
resonant amplitude.

For $r\ge1$, coordinate integration by parts gives the particularly
simple primitive law
\begin{equation}\label{ct:loggamma-primitive}
 \operatorname{FP}\int_0^1\log\Gamma(x)K^{(r)}(x)\,dx
       =-\operatorname{FP}\int_0^1\psi(x)K^{(r-1)}(x)\,dx.
\end{equation}
Here the boundary constant difference is zero, and this can be checked
rather than assumed. For $r=1$, the constant of $\log\Gamma(x)K(x)$
at each endpoint is $-\gamma$. For $r\ge2$, use
\[
 \log\Gamma(x)=-\log x-\gamma x+
     \sum_{j\ge2}\frac{(-1)^j\zeta(j)}j x^j,
 \qquad
 \log\Gamma(1-y)=\gamma y+
     \sum_{j\ge2}\frac{\zeta(j)}j y^j.
\]
The constant at both endpoints of
$\log\Gamma(x)K^{(r-1)}(x)$ is
$-(r-1)!\zeta(r)/r$. Logarithmic terms are excluded by the same
constant-term prescription. This proves \eqref{ct:loggamma-primitive}
directly from the cutoff fundamental theorem.
Combining it with \eqref{ct:zero-index} yields
\[
 \operatorname{FP}\int_0^1\log\Gamma(x)K^{(r)}(x)\,dx
 =\begin{cases}
 -\pi^2/2,&r=1,\\
 2(r-1)!\bigl[\zeta'(r)+H_{r-1}\zeta(r)\bigr],&r\ge2\text{ even},\\
 0,&r\ge3\text{ odd}.
 \end{cases}
\]
For any polynomial $P$ with decomposition \eqref{ct:poly-decomp}, its
log-Gamma moment is therefore a finite sum of these evaluations and
\eqref{ct:loggamma-K}, with the constant contribution
$C\log(2\pi)/2$. This is an explicit primitive counterpart of the
coincident reflected algebra.

The even powers have an especially compact expression involving only
ordinary harmonic numbers:
\begin{equation}\label{ct:loggamma-even}
 \boxed{\displaystyle
 \operatorname{FP}\int_0^1\log\Gamma(x)K(x)^{2h}\,dx
 =\frac{(-1)^h\pi^{2h}}2
       \left[\log(2\pi)-H_{2h}+\frac12H_h\right]
       \quad(h\ge0).}
\end{equation}
Here is an independent proof of the normalization. Initially for
$\Re a>2h-1$,
\[
 V_h(a):=\int_0^1\sin(\pi x)^aK(x)^{2h}\,dx
  =\pi^{2h-1}B\left(\frac{a-2h+1}{2},h+\frac12\right).
\]
At $a=0$, all endpoint exponents are even integers; none is $-1$.
Thus this beta continuation is holomorphic there, and its value and
first derivative equal the coordinate finite parts of $K^{2h}$ and
$K^{2h}\log\sin(\pi x)$. The beta identity \cite{DLMFBeta}, Gamma reflection, and the digamma recurrence \cite{DLMFGamma}
give
\[
 V_h(0)=(-1)^h\pi^{2h},\qquad
 V_h'(0)=(-1)^h\pi^{2h}
       \left[-\log2+H_{2h}-\frac12H_h\right].
\]
Finally average the log-Gamma integral with its reflected copy and use
$\log\Gamma(x)+\log\Gamma(1-x)=\log\pi-\log\sin(\pi x)$.
This proves \eqref{ct:loggamma-even} in the same unit coordinates.
For example,
\[
 \operatorname{FP}\int_0^1\log\Gamma(x)K(x)^2\,dx
   =\frac{\pi^2}{2}\bigl[1-\log(2\pi)\bigr],
\]
\[
 \operatorname{FP}\int_0^1\log\Gamma(x)K(x)^4\,dx
   =\frac{\pi^4}{2}\left[\log(2\pi)-\frac43\right].
\]
The beta/reflection proof uses classical formulas; it is included as an
explicit infinite family and as an independent check on the primitive
normalization, without a claim of global priority.

\subsection{Extension to repeated seams at arbitrary shifts}

The same argument completes the separated theorem by rational partial
fractions. Any finite product of $K^{(r_j)}(x-a_j)$ is rational in
$X=e^{2\pi ix}$, bounded at $X=0,\infty$, and has poles only at the
finitely many $e^{2\pi ia_j}$. At each pole, the principal parts of
$K^{(r)}(x-a)$ form a triangular basis. Subtracting them leaves a
constant. Hence a unique finite representation exists as
\[
 C+\sum_a\sum_{r=0}^{M_a-1}c_{a,r}K^{(r)}(x-a),
\]
where $M_a$ is the pole order at that seam. Multiplication by the ordinary
function $\gamma_m^{(p)}(x)$ and coordinate finite-part integration are
linear operations on its local Laurent-log expansion, so this rational
identity remains valid under precisely the stated normalization. The
terms with $a=0$ are evaluated by \eqref{ct:moment}; the other terms use
the already proved separated Stieltjes--cotangent kernel and its
contact-corrected argument derivatives. No limit through colliding
separated distributions is required.

For completeness, if
$I_m(a)=\operatorname{FP}\int_0^1\gamma_m(x)K(x-a)dx$ at $0<a<1$,
$N=p+r$, and
$h_{m,p}=(-1)^m m!e_{m+1}(p)$ for $p\ge1$ (zero for $p=0$),
the existing one-factor contact law \cite{RepoMain,RepoTriple} gives
\[
 \operatorname{FP}\int_0^1\gamma_m^{(p)}(x)K^{(r)}(x-a)dx
       =(-1)^r\bigl[I_m^{(N)}(a)-h_{m,p}K^{(N)}(a)\bigr].
\]
Here $e_j(p)$ is the elementary symmetric polynomial of degree $j$ in
$1,1/2,\ldots,1/p$, and is zero for $j>p$. To check signs, use
$F_{m,p}=D^pG_m+h_{m,p}\delta^{(p)}$ and
$K^{(N)}(-a)=(-1)^{N+1}K^{(N)}(a)$. Thus arbitrary repetitions and
coincidences within the reflected algebra still reduce to finitely many
ordinary Hurwitz-zeta jets, with trigonometric coefficients at the
remaining distinct shifts.

\subsection{Verification and remaining questions}

The accompanying verifier computes independent endpoint Laurent-log
series from the ordinary Hurwitz shift relation and integrates them
term by term on $(0,1/5)$ and $(4/5,1)$. On the central interval it uses
ordinary high-precision quadrature of Stieltjes/polygamma and cotangent
functions. It compares these values against finite coefficients of
\eqref{ct:R}, and separately integrates powers $K^q$ without applying
the derivative-polynomial reduction. Exact symbolic assertions verify
the power reduction recurrence. These are independent floating-point
diagnostics, not interval enclosures; the proof above establishes the
all-index result.

The natural next questions are to classify which products with two
unreflected Stieltjes factors and a reflected polynomial still close
without a Tornheim jet, to derive a conductor-compatible version of the
repeated-seam partial fractions, and to formalize the local coefficient
completion independently of distributional products. The present
results do not imply an arithmetic evaluation of the generic cubic
Tornheim derivative or either fixed-weight Gaussian candidate.
